\begin{thebibliography}{99}

\bibitem{Chen1966} J. R. Chen,
On the representation of a large even integer as the sum of a prime and
the product of at most two primes (in Chinese),
\emph{Kexue Tongbao} \textbf{17} (1966), 385--386.

\bibitem{Chen1973} J. R. Chen,
On the representation of a large even integer as the sum of a prime and
the product of at most two primes,
\emph{Scientia Sinica} \textbf{16} (1973), 157--176.

\bibitem{Wu2004} J. Wu,
Chen's double sieve, Goldbach's conjecture and the twin prime problem,
\emph{Acta Arithmetica} \textbf{114} (2004), no.~3, 215--273.

\bibitem{Wu2008} J. Wu,
Chen's double sieve, Goldbach's conjecture and the twin prime problem, 2,
\emph{Acta Arithmetica} \textbf{131} (2008), no.~4, 367--387.
DOI: \href{https://doi.org/10.4064/aa131-4-5}{10.4064/aa131-4-5}.

\bibitem{Lichtman} J. D. Lichtman,
Primes in arithmetic progressions to large moduli, and Goldbach beyond
the square-root barrier, with an appendix with S. Drappeau,
\href{https://arxiv.org/abs/2309.08522v1}{arXiv:2309.08522v1}, 2023.

\bibitem{LiLiu2026} J. Li and J. Liu,
Theorem $(1+1.9)$ on the Goldbach Conjecture,
\href{https://arxiv.org/abs/2606.05224v2}{arXiv:2606.05224v2}, 2026.

\bibitem[Z1]{Z1} R. Zhang,
Shifted primes in arithmetic progressions to large moduli with linear sieve remainder weights,
preprint, Zenodo, 2026.
DOI: \href{https://doi.org/10.5281/zenodo.23615525}{10.5281/zenodo.23615525}.

\bibitem{ref1} J.-M. Deshouillers and H. Iwaniec,
Kloosterman sums and Fourier coefficients of cusp forms,
\emph{Inventiones Mathematicae} \textbf{70} (1982), 219--288.
DOI: \href{https://doi.org/10.1007/BF01390728}{10.1007/BF01390728}.

\bibitem{ref2} E. Bombieri, J. B. Friedlander and H. Iwaniec,
Primes in arithmetic progressions to large moduli,
\emph{Acta Mathematica} \textbf{156} (1986), 203--251.
DOI: \href{https://doi.org/10.1007/BF02399204}{10.1007/BF02399204}.

\bibitem{ref3} S. Drappeau,
Sums of Kloosterman sums in arithmetic progressions, and the error term
in the dispersion method,
\emph{Proceedings of the London Mathematical Society} \textbf{114}
(2017), no.~4, 684--732.
DOI: \href{https://doi.org/10.1112/plms.12022}{10.1112/plms.12022};
\href{https://arxiv.org/abs/1504.05549v4}{arXiv:1504.05549v4}.

\bibitem{ref4} H. Iwaniec,
Rosser's sieve,
\emph{Acta Arithmetica} \textbf{36} (1980), 171--202.

\bibitem{ref5} H. Iwaniec,
A new form of the error term in the linear sieve,
\emph{Acta Arithmetica} \textbf{37} (1980), 307--320.
DOI: \href{https://doi.org/10.4064/aa-37-1-307-320}{10.4064/aa-37-1-307-320}.

\bibitem{ref6} D. R. Heath-Brown,
Prime numbers in short intervals and a generalized Vaughan identity,
\emph{Canadian Journal of Mathematics} \textbf{34} (1982), 1365--1377.
DOI: \href{https://doi.org/10.4153/CJM-1982-095-9}{10.4153/CJM-1982-095-9}.

\bibitem{ref7} H. H. Kim,
Functoriality for the exterior square of $GL_4$ and the symmetric fourth
of $GL_2$,
\emph{Journal of the American Mathematical Society} \textbf{16}
(2003), 139--183, with Appendix~1 by D. Ramakrishnan and Appendix~2
by H. H. Kim and P. Sarnak.
DOI: \href{https://doi.org/10.1090/S0894-0347-02-00410-1}{10.1090/S0894-0347-02-00410-1}.

\bibitem{ref11} H. Iwaniec and E. Kowalski,
\emph{Analytic Number Theory}, American Mathematical Society Colloquium
Publications \textbf{53}, American Mathematical Society, Providence, RI, 2004.

\bibitem{Shiu1980} P. Shiu,
A Brun--Titchmarsh theorem for multiplicative functions,
\emph{Journal f\"ur die reine und angewandte Mathematik} \textbf{313}
(1980), 161--170.

\bibitem{Bombieri1965} E. Bombieri,
On the large sieve,
\emph{Mathematika} \textbf{12} (1965), 201--225.

\bibitem{Vinogradov1965} A. I. Vinogradov,
The density hypothesis for Dirichlet $L$-series (in Russian),
\emph{Izvestiya Akademii Nauk SSSR. Seriya Matematicheskaya}
\textbf{29} (1965), 903--934.

\bibitem{Siegel1935} C. L. Siegel,
\"Uber die Classenzahl quadratischer Zahlk\"orper,
\emph{Acta Arithmetica} \textbf{1} (1935), 83--86.

\bibitem{Walfisz1936} A. Walfisz,
Zur additiven Zahlentheorie.~II,
\emph{Mathematische Zeitschrift} \textbf{40} (1936), 592--607.
DOI: \href{https://doi.org/10.1007/BF01218882}{10.1007/BF01218882}.

\bibitem{Mertens1874} F. Mertens,
Ein Beitrag zur analytischen Zahlentheorie,
\emph{Journal f\"ur die reine und angewandte Mathematik} \textbf{78}
(1874), 46--62.

\bibitem{Buchstab1937} A. A. Buchstab,
Asymptotische Absch\"atzungen einer allgemeinen zahlentheoretischen Funktion,
\emph{Matematicheskii Sbornik (N.S.)} \textbf{2(44)} (1937),
no.~6, 1239--1246.

\bibitem{RosserSchoenfeld1962} J. B. Rosser and L. Schoenfeld,
Approximate formulas for some functions of prime numbers,
\emph{Illinois Journal of Mathematics} \textbf{6} (1962), 64--94.
DOI: \href{https://doi.org/10.1215/ijm/1255631807}{10.1215/ijm/1255631807}.

\bibitem{Johansson2017} F. Johansson,
Arb: Efficient arbitrary-precision midpoint-radius interval arithmetic,
\emph{IEEE Transactions on Computers} \textbf{66} (2017), no.~8,
1281--1292.
DOI: \href{https://doi.org/10.1109/TC.2017.2690633}{10.1109/TC.2017.2690633}.

\end{thebibliography}
